% !TEX program = lualatex
% ---------------------------------------------------------------------------------------------
%  Readback note 15: A crude bound for a class number, and divisibility of indices of submodules
%
%  Build:  python3 tools/rb.py note <slug>
%
%  One unit of a Lean artifact, rewritten as ordinary mathematics. The prose is hand-written. Lean names,
%  line ranges and counts come from data.tex; quoted Lean statements are copied from the artifact by the build.
%
%  Rules of the text: standard terminology; every symbol is explained where it first appears and listed again
%  in the notation at the end; no letter has two meanings; a proof is complete or says what it takes for granted;
%  Lean names appear only in the LEAN lines and in the concordance.
%
%  The signature rule (the build enforces it): a function, relation or predicate that the artifact defines is
%  written only through its notation macro, which the build makes from the spec (`notation`) with one slot for
%  each explicit argument of the Lean definition. No argument is fixed and dropped, also not "for brevity" and
%  not for one section. Every argument and every hypothesis of a stated declaration stands inside its
%  statement, or in a Setting that the statement cites by number; only implicit and instance arguments (the
%  ambient objects) may stand in the conventions.
% ---------------------------------------------------------------------------------------------
\documentclass[11pt,a4paper]{article}
\ifdefined\pdfvariable\pdfvariable suppressoptionalinfo \numexpr 1+2+4+8\relax\fi
\usepackage[a4paper,hmargin=35mm,top=22mm,bottom=27mm,footskip=12mm,nohead]{geometry}

\usepackage{fontspec}
\directlua{
  luaotfload.add_fallback("gfallback", {"STIXTwoMath-Regular.otf:mode=node;", "DejaVuSans.ttf:mode=node;"})
  luaotfload.add_fallback("gmonofallback", {"DejaVuSansMono.ttf:mode=node;", "STIXTwoMath-Regular.otf:mode=node;"})
}
\setmainfont{NewCM10}[Extension=.otf,UprightFont=*-Book,ItalicFont=*-BookItalic,
  BoldFont=*-Bold,BoldItalicFont=*-BoldItalic,SmallCapsFont=*-Book,SmallCapsFeatures={Letters=SmallCaps},
  RawFeature={fallback=gfallback}]
\setsansfont{NewCMSans10}[Extension=.otf,UprightFont=*-Book,ItalicFont=*-BookOblique,
  BoldFont=*-Bold,BoldItalicFont=*-BoldOblique,RawFeature={fallback=gfallback}]
\setmonofont{JetBrainsMono}[Extension=.otf,UprightFont=*-Regular,ItalicFont=*-Italic,
  BoldFont=*-SemiBold,BoldItalicFont=*-SemiBoldItalic,Scale=MatchLowercase,
  RawFeature={-calt;-liga;fallback=gmonofallback}]
\usepackage{amsmath}
\usepackage{amsthm}
\usepackage[math-style=TeX]{unicode-math}
\setmathfont{NewCMMath-Book.otf}
\emergencystretch=2.5em\relax

\usepackage{xcolor}
\usepackage{tikz}
\usetikzlibrary{calc,arrows.meta,positioning,backgrounds}
\usepackage[protrusion=true,expansion=true]{microtype}
\usepackage{enumitem}
\usepackage{needspace}
\usepackage{booktabs}
\usepackage{array}
\usepackage{longtable}
\usepackage{fancyhdr}
\usepackage{lastpage}
\usepackage{multicol}
\usepackage[hidelinks,pdfusetitle]{hyperref}

\definecolor{ink}{HTML}{0F172A}
\definecolor{inksoft}{HTML}{475569}
\definecolor{inkfaint}{HTML}{94A3B8}
\definecolor{rulegrey}{HTML}{CBD5E1}
\definecolor{arith}{HTML}{C2410C}
\definecolor{est}{HTML}{7C3AED}
\definecolor{ana}{HTML}{0891B2}
\definecolor{ok}{HTML}{047857}
\definecolor{leancol}{HTML}{3730A3}

\title{A crude bound for a class number, and divisibility of indices of submodules}
\author{grh-compact}
\date{}

% ------------------------------------------------------------------------------------- macros
\newcommand{\leanfont}{\ttfamily\fontsize{8.6}{10.5}\selectfont\color{leancol}}
\newcommand{\lean}[1]{\csname LN@#1\endcsname}
\newcommand{\fillin}[1]{{\color{arith}\itshape[#1]}}  % a place of the template that is not written yet; the build reports every one that is left
\newcommand{\leanquote}[1]{\csname LQ@#1\endcsname}
% generated by rb.py note; do not edit
\newcommand{\ConcDecls}{20}
\newcommand{\Pin}{72fead10}
\newcommand{\PinText}{that of commit {\leanfont 72fead10}}
\newcommand{\ArtifactName}{grh-compact}
\expandafter\newcommand\csname LN@cardclassesle\endcsname{\href{https://github.com/lsernau/grh-lean/blob/72fead102f663d5c7a407964a13dc0cf9d977207/grh-compact/src/New/Core/ClassCountBound.lean\#L14-L37}{{\leanfont Space.\allowbreak{}card\_classes\_le}}}
\expandafter\newcommand\csname LQ@cardclassesle\endcsname{\leanquoteline{0}{lemma Space.card\_classes\_le \{E\} [Field E] [CMType E] \{N\}}\leanquoteline{4}{[NeZero N] (h : ↑N ∈ differentIdeal ℤ (𝓞 E)) \{n\} [Fintype n] : .card (StdClasses}\leanquoteline{4}{(Kf E n)) ≤ (1 + 2 * shortBasisConst E n N \textasciicircum{} 3) \textasciicircum{} (4 * numInfPlaces E * card n \textasciicircum{} 2)}}
\expandafter\newcommand\csname LN@IndexDvd\endcsname{\href{https://github.com/lsernau/grh-lean/blob/72fead102f663d5c7a407964a13dc0cf9d977207/grh-compact/src/New/Core/ModuleIndex.lean\#L20-L20}{{\leanfont IndexDvd}}}
\expandafter\newcommand\csname LQ@IndexDvd\endcsname{\leanquoteline{0}{def IndexDvd := ∀ n (Φ : Fin n → V), (∀ k, Φ k ∈ M) → ∃ m, det m ∣ z ∧ ∀ i, ∑ k, m i k • Φ k ∈ N}}
\expandafter\newcommand\csname LN@comap\endcsname{\href{https://github.com/lsernau/grh-lean/blob/72fead102f663d5c7a407964a13dc0cf9d977207/grh-compact/src/New/Core/ModuleIndex.lean\#L24-L25}{{\leanfont IndexDvd.\allowbreak{}comap}}}
\expandafter\newcommand\csname LN@ofle\endcsname{\href{https://github.com/lsernau/grh-lean/blob/72fead102f663d5c7a407964a13dc0cf9d977207/grh-compact/src/New/Core/ModuleIndex.lean\#L28-L28}{{\leanfont IndexDvd.\allowbreak{}of\_le}}}
\expandafter\newcommand\csname LN@trans\endcsname{\href{https://github.com/lsernau/grh-lean/blob/72fead102f663d5c7a407964a13dc0cf9d977207/grh-compact/src/New/Core/ModuleIndex.lean\#L31-L33}{{\leanfont IndexDvd.\allowbreak{}trans}}}
\expandafter\newcommand\csname LN@map\endcsname{\href{https://github.com/lsernau/grh-lean/blob/72fead102f663d5c7a407964a13dc0cf9d977207/grh-compact/src/New/Core/ModuleIndex.lean\#L37-L41}{{\leanfont IndexDvd.\allowbreak{}map}}}
\expandafter\newcommand\csname LN@sup\endcsname{\href{https://github.com/lsernau/grh-lean/blob/72fead102f663d5c7a407964a13dc0cf9d977207/grh-compact/src/New/Core/ModuleIndex.lean\#L46-L48}{{\leanfont IndexDvd.\allowbreak{}sup}}}
\expandafter\newcommand\csname LN@comp\endcsname{\href{https://github.com/lsernau/grh-lean/blob/72fead102f663d5c7a407964a13dc0cf9d977207/grh-compact/src/New/Core/ModuleIndex.lean\#L50-L51}{{\leanfont IndexDvd.\allowbreak{}comp}}}
\expandafter\newcommand\csname LN@HasType\endcsname{\href{https://github.com/lsernau/grh-lean/blob/72fead102f663d5c7a407964a13dc0cf9d977207/grh-compact/src/New/Core/ModuleIndex.lean\#L55-L56}{{\leanfont HasType}}}
\expandafter\newcommand\csname LQ@HasType\endcsname{\leanquoteline{0}{def HasType p M n (y : End O V) (c : ℝ) :=}\leanquoteline{2}{M.map y ≤ M ∧ ∃ l : ℕ, l ≤ c ∧ IndexDvd M (M.map y ⊔ p \textasciicircum{} n • M) (p \textasciicircum{} l)}}
\expandafter\newcommand\csname LN@prod\endcsname{\href{https://github.com/lsernau/grh-lean/blob/72fead102f663d5c7a407964a13dc0cf9d977207/grh-compact/src/New/Core/ModuleIndex.lean\#L60-L64}{{\leanfont HasType.\allowbreak{}prod}}}
\expandafter\newcommand\csname LN@ofdetdvd\endcsname{\href{https://github.com/lsernau/grh-lean/blob/72fead102f663d5c7a407964a13dc0cf9d977207/grh-compact/src/New/Core/ModuleIndex.lean\#L68-L73}{{\leanfont IndexDvd.\allowbreak{}of\_det\_dvd}}}
\expandafter\newcommand\csname LQ@ofdetdvd\endcsname{\leanquoteline{0}{lemma IndexDvd.of\_det\_dvd [IsDomain O] [ValuationRing O] [Fintype ι] \{A : Matrix ι ι O\} \{N\}}\leanquoteline{4}{(h : A.det ∣ z) (e : .map A.mulVecLin ⊤ ≤ N) : IndexDvd ⊤ N z}}
\expandafter\newcommand\csname LN@numInf\endcsname{\href{https://github.com/lsernau/grh-lean/blob/72fead102f663d5c7a407964a13dc0cf9d977207/grh-compact/src/Lib/General/Analysis/ZpowSum.lean\#L32-L32}{{\leanfont numInfPlaces}}}
\expandafter\newcommand\csname LN@sbc\endcsname{\href{https://github.com/lsernau/grh-lean/blob/72fead102f663d5c7a407964a13dc0cf9d977207/grh-compact/src/Lib/Adelic/Transference/MassBasic.lean\#L73-L73}{{\leanfont Space.\allowbreak{}shortBasisConst}}}
\expandafter\newcommand\csname LN@Kf\endcsname{\href{https://github.com/lsernau/grh-lean/blob/72fead102f663d5c7a407964a13dc0cf9d977207/grh-compact/src/Lib/Adelic/Unitary/SpaceDefs.lean\#L109-L109}{{\leanfont Kf}}}
\expandafter\newcommand\csname LN@StdClasses\endcsname{\href{https://github.com/lsernau/grh-lean/blob/72fead102f663d5c7a407964a13dc0cf9d977207/grh-compact/src/Lib/Adelic/CuspForms/ClassReps.lean\#L93-L93}{{\leanfont StdClasses}}}
\expandafter\newcommand\csname LN@shortbasis\endcsname{\href{https://github.com/lsernau/grh-lean/blob/72fead102f663d5c7a407964a13dc0cf9d977207/grh-compact/src/Lib/Adelic/Transference/MassBasic.lean\#L101-L107}{{\leanfont Space.\allowbreak{}exists\_short\_basis\_stdSpace}}}
\expandafter\newcommand\csname LN@ratnear\endcsname{\href{https://github.com/lsernau/grh-lean/blob/72fead102f663d5c7a407964a13dc0cf9d977207/grh-compact/src/Lib/Adelic/Transference/MassBasic.lean\#L116-L124}{{\leanfont exists\_rat\_near\_of\_mem\_differentIdeal}}}
\expandafter\newcommand\csname LN@enc\endcsname{\href{https://github.com/lsernau/grh-lean/blob/72fead102f663d5c7a407964a13dc0cf9d977207/grh-compact/src/New/Filling/BoxPartition.lean\#L287-L292}{{\leanfont Filling.\allowbreak{}encard\_integral\_le}}}
\expandafter\newcommand\csname LN@hermite\endcsname{\href{https://github.com/lsernau/grh-lean/blob/72fead102f663d5c7a407964a13dc0cf9d977207/grh-compact/src/Lib/Adelic/Witt/FrameAction.lean\#L58-L69}{{\leanfont IsValuationPair.\allowbreak{}hermite}}}
\expandafter\newcommand\csname LN@valtop\endcsname{\href{https://github.com/lsernau/grh-lean/blob/72fead102f663d5c7a407964a13dc0cf9d977207/grh-compact/src/Lib/Adelic/Witt/FrameAction.lean\#L50-L51}{{\leanfont isValuationPair\_top}}}
\newcommand{\ConcordanceRows}{\crow{cardclassesle}{\href{https://github.com/lsernau/grh-lean/blob/72fead102f663d5c7a407964a13dc0cf9d977207/grh-compact/src/New/Core/ClassCountBound.lean\#L14-L37}{{\leanfont Space.\allowbreak{}card\_classes\_le}}}{ClassCountBound.lean}{14--37}
\crow{IndexDvd}{\href{https://github.com/lsernau/grh-lean/blob/72fead102f663d5c7a407964a13dc0cf9d977207/grh-compact/src/New/Core/ModuleIndex.lean\#L20-L20}{{\leanfont IndexDvd}}}{ModuleIndex.lean}{20--20}
\crow{comap}{\href{https://github.com/lsernau/grh-lean/blob/72fead102f663d5c7a407964a13dc0cf9d977207/grh-compact/src/New/Core/ModuleIndex.lean\#L24-L25}{{\leanfont IndexDvd.\allowbreak{}comap}}}{ModuleIndex.lean}{24--25}
\crow{ofle}{\href{https://github.com/lsernau/grh-lean/blob/72fead102f663d5c7a407964a13dc0cf9d977207/grh-compact/src/New/Core/ModuleIndex.lean\#L28-L28}{{\leanfont IndexDvd.\allowbreak{}of\_le}}}{ModuleIndex.lean}{28--28}
\crow{trans}{\href{https://github.com/lsernau/grh-lean/blob/72fead102f663d5c7a407964a13dc0cf9d977207/grh-compact/src/New/Core/ModuleIndex.lean\#L31-L33}{{\leanfont IndexDvd.\allowbreak{}trans}}}{ModuleIndex.lean}{31--33}
\crow{map}{\href{https://github.com/lsernau/grh-lean/blob/72fead102f663d5c7a407964a13dc0cf9d977207/grh-compact/src/New/Core/ModuleIndex.lean\#L37-L41}{{\leanfont IndexDvd.\allowbreak{}map}}}{ModuleIndex.lean}{37--41}
\crow{sup}{\href{https://github.com/lsernau/grh-lean/blob/72fead102f663d5c7a407964a13dc0cf9d977207/grh-compact/src/New/Core/ModuleIndex.lean\#L46-L48}{{\leanfont IndexDvd.\allowbreak{}sup}}}{ModuleIndex.lean}{46--48}
\crow{comp}{\href{https://github.com/lsernau/grh-lean/blob/72fead102f663d5c7a407964a13dc0cf9d977207/grh-compact/src/New/Core/ModuleIndex.lean\#L50-L51}{{\leanfont IndexDvd.\allowbreak{}comp}}}{ModuleIndex.lean}{50--51}
\crow{HasType}{\href{https://github.com/lsernau/grh-lean/blob/72fead102f663d5c7a407964a13dc0cf9d977207/grh-compact/src/New/Core/ModuleIndex.lean\#L55-L56}{{\leanfont HasType}}}{ModuleIndex.lean}{55--56}
\crow{prod}{\href{https://github.com/lsernau/grh-lean/blob/72fead102f663d5c7a407964a13dc0cf9d977207/grh-compact/src/New/Core/ModuleIndex.lean\#L60-L64}{{\leanfont HasType.\allowbreak{}prod}}}{ModuleIndex.lean}{60--64}
\crow{ofdetdvd}{\href{https://github.com/lsernau/grh-lean/blob/72fead102f663d5c7a407964a13dc0cf9d977207/grh-compact/src/New/Core/ModuleIndex.lean\#L68-L73}{{\leanfont IndexDvd.\allowbreak{}of\_det\_dvd}}}{ModuleIndex.lean}{68--73}
\crow{numInf}{\href{https://github.com/lsernau/grh-lean/blob/72fead102f663d5c7a407964a13dc0cf9d977207/grh-compact/src/Lib/General/Analysis/ZpowSum.lean\#L32-L32}{{\leanfont numInfPlaces}}}{ZpowSum.lean}{32--32}
\crow{sbc}{\href{https://github.com/lsernau/grh-lean/blob/72fead102f663d5c7a407964a13dc0cf9d977207/grh-compact/src/Lib/Adelic/Transference/MassBasic.lean\#L73-L73}{{\leanfont Space.\allowbreak{}shortBasisConst}}}{MassBasic.lean}{73--73}
\crow{Kf}{\href{https://github.com/lsernau/grh-lean/blob/72fead102f663d5c7a407964a13dc0cf9d977207/grh-compact/src/Lib/Adelic/Unitary/SpaceDefs.lean\#L109-L109}{{\leanfont Kf}}}{SpaceDefs.lean}{109--109}
\crow{StdClasses}{\href{https://github.com/lsernau/grh-lean/blob/72fead102f663d5c7a407964a13dc0cf9d977207/grh-compact/src/Lib/Adelic/CuspForms/ClassReps.lean\#L93-L93}{{\leanfont StdClasses}}}{ClassReps.lean}{93--93}
\crow{shortbasis}{\href{https://github.com/lsernau/grh-lean/blob/72fead102f663d5c7a407964a13dc0cf9d977207/grh-compact/src/Lib/Adelic/Transference/MassBasic.lean\#L101-L107}{{\leanfont Space.\allowbreak{}exists\_short\_basis\_stdSpace}}}{MassBasic.lean}{101--107}
\crow{ratnear}{\href{https://github.com/lsernau/grh-lean/blob/72fead102f663d5c7a407964a13dc0cf9d977207/grh-compact/src/Lib/Adelic/Transference/MassBasic.lean\#L116-L124}{{\leanfont exists\_rat\_near\_of\_mem\_differentIdeal}}}{MassBasic.lean}{116--124}
\crow{enc}{\href{https://github.com/lsernau/grh-lean/blob/72fead102f663d5c7a407964a13dc0cf9d977207/grh-compact/src/New/Filling/BoxPartition.lean\#L287-L292}{{\leanfont Filling.\allowbreak{}encard\_integral\_le}}}{BoxPartition.lean}{287--292}
\crow{hermite}{\href{https://github.com/lsernau/grh-lean/blob/72fead102f663d5c7a407964a13dc0cf9d977207/grh-compact/src/Lib/Adelic/Witt/FrameAction.lean\#L58-L69}{{\leanfont IsValuationPair.\allowbreak{}hermite}}}{FrameAction.lean}{58--69}
\crow{valtop}{\href{https://github.com/lsernau/grh-lean/blob/72fead102f663d5c7a407964a13dc0cf9d977207/grh-compact/src/Lib/Adelic/Witt/FrameAction.lean\#L50-L51}{{\leanfont isValuationPair\_top}}}{FrameAction.lean}{50--51}}
% the notations of the defined functions: one slot for each explicit argument of the Lean declaration
\newcommand{\IndexDvd}[3]{[#1\,{:}\,#2]\mathrel{\big|}#3}
\newcommand{\HasType}[5]{$#4$ has type at most~$#5$ at~$#1^{#3}$ on~$#2$}
\newcommand{\numInf}[1]{f(#1)}
\newcommand{\sbc}[3]{c_{\mathrm{sb}}(#1,#2,#3)}
\newcommand{\Kf}[2]{\mathrm{U}_{#2}(\widehat{\mathfrak{o}}_{#1})}
\newcommand{\StdClasses}[1]{\mathrm{Cl}(#1)}
\newcommand{\UnitMembers}{11}
\newcommand{\UnitInNote}{11}
\newcommand{\UnitProofs}{5}
\newcommand{\UnitSteps}{22}
\newcommand{\UnitAllSteps}{24}


\linespread{1.08}                       % 11 pt on 14.7 pt
\setlength{\parindent}{0pt}
\setlength{\parskip}{5.5pt plus 1.5pt minus 0.5pt}
\setlength{\abovedisplayskip}{7.5pt plus 2pt minus 3pt}
\setlength{\belowdisplayskip}{7.5pt plus 2pt minus 3pt}
\setlength{\abovedisplayshortskip}{3pt plus 2pt}
\setlength{\belowdisplayshortskip}{6pt plus 2pt minus 2pt}
\emergencystretch=1.5em
\clubpenalty=10000 \widowpenalty=10000 \displaywidowpenalty=10000
\relpenalty=7000 \binoppenalty=9000
\raggedbottom
\pagestyle{fancy}\fancyhf{}\renewcommand{\headrulewidth}{0pt}
\fancyfoot[C]{\sffamily\fontsize{7.6}{9}\selectfont\color{inkfaint}\thepage\,/\,\pageref*{LastPage}}
\setlist{topsep=2pt,partopsep=0pt,parsep=0pt,itemsep=4pt,leftmargin=2.1em,labelsep=0.55em}

\newcommand{\Q}{\mathbb{Q}}
\newcommand{\Z}{\mathbb{Z}}
\newcommand{\F}{\mathbb{F}}
\newcommand{\len}{\operatorname{len}}
\newcommand{\cok}{\operatorname{cok}}
\newcommand{\End}{\operatorname{End}}
\newcommand{\Hom}{\operatorname{Hom}}
\newcommand{\Gal}{\operatorname{Gal}}
\newcommand{\diag}{\operatorname{diag}}
\newcommand{\kw}[1]{\emph{#1}}
\newcommand{\m}[1]{\mbox{$#1$}}        % a short formula that must not be broken across lines
\newcommand{\head}[1]{\par\addvspace{17pt plus 4pt minus 2pt}\Needspace{6\baselineskip}%
  \noindent{\fontsize{13.2}{16}\selectfont\bfseries\color{ink}#1}\par\nopagebreak\vspace{-1pt}}
\newcommand{\sect}[2]{\refstepcounter{section}\label{#2}\par\addvspace{19pt plus 4pt minus 2pt}\Needspace{7\baselineskip}%
  \noindent\llap{{\fontsize{13.2}{16}\selectfont\bfseries\color{inkfaint}\thesection}\hspace{9pt}}%
  {\fontsize{13.2}{16}\selectfont\bfseries\color{ink}#1}\par\nopagebreak\vspace{-1pt}}
\newcommand{\subhead}[1]{\par\addvspace{9pt plus 2pt}\Needspace{4\baselineskip}\noindent\textbf{#1}\enspace\ignorespaces}
% where a statement lives in the Lean archive
\newcommand{\leantag}{{\sffamily\bfseries\fontsize{6.3}{8}\selectfont\color{leancol}\addfontfeatures{LetterSpace=7}LEAN}}
\newcommand{\sep}{\unskip\nobreak\hspace{5pt}{\color{inkfaint}\textperiodcentered}\hspace{5pt plus 2pt}\ignorespaces}
\newcommand{\leanline}[1]{\par\nopagebreak\vspace{-2.5pt}{\fontsize{8.9}{12.4}\selectfont\color{inksoft}\raggedright
  \noindent\llap{\leantag\hspace{9pt}}#1\par}\vspace{1pt}}
% a displayed statement in a tinted box: #1 colour, #2 text
\newcommand{\statement}[2]{\par\addvspace{7pt}\begingroup\setlength{\fboxsep}{8.5pt}\setlength{\fboxrule}{0.6pt}%
  \noindent\fcolorbox{#1!30!white}{#1!3!white}{\parbox{\dimexpr\linewidth-2\fboxsep-2\fboxrule\relax}{#2}}%
  \endgroup\par\addvspace{7pt}}
\newcommand{\nt}[2]{\par\noindent\hangindent=1.1em\hangafter=1 #1\enspace{\color{inksoft}#2}\par\vspace{1pt}}
\newcommand{\legdot}[1]{\tikz[baseline=-0.62ex]\fill[#1] (0,0) circle (2.5pt);\hspace{3pt}}
\newcommand{\crow}[4]{\csname CS@#1\endcsname & #2 & \textcolor{inksoft}{#3} & \textcolor{inksoft}{#4}\\}
\newcommand{\leanquoteline}[2]{\par\noindent\hangindent=\dimexpr#1\fontdimen2\font+2.2em\relax\hangafter=1
  \hspace*{#1\fontdimen2\font}#2}
\newcommand{\quoted}[1]{\par\vspace{3pt}{\ttfamily\fontsize{7.5}{10.4}\selectfont\color{leancol}\raggedright
  \setlength{\parskip}{0pt}\leanquote{#1}\par}\vspace{4pt}}
\newcommand{\bx}[3]{{\sffamily\bfseries\fontsize{7.3}{9.2}\selectfont\color{#1!82!black}#2}\\[1.5pt]#3}
\newcommand{\cs}[2]{\expandafter\newcommand\csname CS@#1\endcsname{#2}}

% theorems
\newtheoremstyle{plainx}{7pt plus 2pt}{7pt plus 2pt}{\itshape}{}{\bfseries}{.}{0.5em}{}
\newtheoremstyle{uprightx}{7pt plus 2pt}{7pt plus 2pt}{}{}{\bfseries}{.}{0.5em}{}
\theoremstyle{plainx}
\newtheorem{theorem}{Theorem}[section]
\newtheorem{lemma}[theorem]{Lemma}
\newtheorem{proposition}[theorem]{Proposition}
\newtheorem{corollary}[theorem]{Corollary}
\theoremstyle{uprightx}
\newtheorem{definition}[theorem]{Definition}
\newtheorem{setting}[theorem]{Setting}          % hypotheses shared by several statements; each of them cites it by number
\newtheorem{remark}[theorem]{Remark}
\newtheorem{example}[theorem]{Example}
\renewcommand{\qedsymbol}{{\setlength{\fboxsep}{0pt}\setlength{\fboxrule}{0.5pt}\fbox{\rule{0pt}{0.46em}\hspace{0.46em}}}}
% what holds throughout the note: the ambient objects only (implicit and instance arguments), never a hypothesis
\newenvironment{conventions}{\par\addvspace{5pt}\noindent\textbf{Conventions.}\enspace\ignorespaces}{\par\addvspace{3pt}}
\makeatletter
\renewenvironment{proof}[1][Proof]{\par\pushQED{\qed}\normalfont\topsep0pt\partopsep0pt
  \trivlist\item[\hskip\labelsep\itshape #1.]\ignorespaces}{\popQED\endtrivlist\@endpefalse}
\makeatother

% where each declaration of the concordance is treated: \cs{key}{place}, the place naming a \label or a lettered item
\cs{cardclassesle}{Theorem 1}
\cs{ofdetdvd}{Theorem 2}
\cs{IndexDvd}{Definition~\ref{d:index}}
\cs{HasType}{Definition~\ref{d:type}}
\cs{Kf}{Definition~\ref{d:classes}}
\cs{StdClasses}{Definition~\ref{d:classes}}
\cs{numInf}{Definition~\ref{d:classes}}
\cs{sbc}{Definition~\ref{d:classes}}
\cs{shortbasis}{Input 1}
\cs{ratnear}{Input 2}
\cs{enc}{Input 3}
\cs{hermite}{Input 4}
\cs{valtop}{Input 4}
\cs{ofle}{Lemma~\ref{l:incl}}
\cs{comap}{Lemma~\ref{l:comap}}
\cs{trans}{Lemma~\ref{l:trans}}
\cs{map}{Lemma~\ref{l:map}}
\cs{sup}{Lemma~\ref{l:sup}}
\cs{comp}{Lemma~\ref{l:comp}}
\cs{prod}{Proposition~\ref{p:prod}}
\newcommand{\cO}{\mathcal{O}}
\newcommand{\cJ}{\mathcal{J}}
\newcommand{\cS}{\mathcal{S}}
\newcommand{\Af}{\mathbb{A}_{E,f}}
\newcommand{\Oh}{\widehat{\mathcal{O}}_{E}}
\newcommand{\UJ}{\mathrm{U}_{\mathcal{J}}}
\newcommand{\GLJ}{\mathrm{GL}_{\mathcal{J}}}
\newcommand{\MJ}{\mathrm{M}_{\mathcal{J}}}
\newcommand{\fc}{\mathfrak{c}}
\newcommand{\CC}{\mathbb{C}}

\begin{document}
\color{ink}

{\raggedright
{\sffamily\bfseries\fontsize{7.2}{9}\selectfont\color{leancol}\addfontfeatures{LetterSpace=8}READBACK NOTE 15\par}
\vspace{5pt}
{\fontsize{21.5}{25.5}\selectfont\bfseries A crude bound for a class number, and divisibility of indices of submodules\par}
\vspace{5pt}
{\fontsize{12.4}{16}\selectfont\itshape\color{inksoft} One unit of the Lean artifact \ArtifactName, read back into
mathematics\par}}
\vspace{8pt}
{\color{rulegrey}\hrule height 0.6pt}
\vspace{7pt}
\thispagestyle{empty}

\textbf{Summary.} The unit holds two pieces of bookkeeping that do not depend on each other. The first is an
explicit, crude bound for the number of classes in the genus of the standard lattice of a definite hermitian
space over a CM field: each class is recorded by two matrices of algebraic integers of bounded size, and such
matrices are counted. The second is a way to say, over an arbitrary commutative ring, that the index of one
submodule in another divides a given element: every finite family of the first is moved into the second by a
square matrix whose determinant divides that element. The unit proves the calculus of this notion and a
criterion by determinants over a valuation domain, which rests on the Hermite normal form. Both pieces serve
Theorem~A of the write-up (short forms): the class number enters the sizes of the columns, and the index language
is the one in which the losses of divisibility at single primes are stated.

\statement{leancol}{\fontsize{9.6}{13.2}\selectfont\textbf{What this note is.} The artifact \ArtifactName{} is a
body of Lean code: the compact edition of a formal proof, due to L.~Alpöge with the AI system Claude, of the
Riemann hypothesis and of its generalisation to Dirichlet $L$-functions. That headline is a claim. The authors
report that every file compiles with Lean v4.35.0 against a pinned version of the library Mathlib and that Lean's
kernel accepted the proofs; an independent build is recorded in the handout's record of checks when it reports.
No mathematician has read the argument through, and nothing here says that the Riemann hypothesis is proved. This
note treats one unit of the artifact: all \UnitInNote{} of its \UnitMembers{} declarations. Their proofs are given
in full, \emph{given} four results of other units, which are stated as Inputs~1 to~4 and not proved here: two
lemmas of reduction theory, a count of bounded integral vectors, and the Hermite normal form. Every statement of
the unit and of the Inputs was compared with its Lean text, argument by argument. The Lean text read is
\PinText; Appendix~A is the concordance of \ConcDecls{} declarations.}

\sect{The statements}{s:intro}
The notions are defined in §\ref{s:setting}. In short: for a CM field~$E$ and a finite set~$\cJ$,
\m{\StdClasses{\Kf{E}{\cJ}}} is the set of classes in the genus of the standard lattice of the hermitian space
$E^{\cJ}$ with the form \m{\sum_{i}\bar v_i v_i}; \m{\numInf{E}} is the number of infinite places of~$E$; and
\m{\sbc{E}{\cJ}{M}=6M(\numInf{E}\cdot\#\cJ)^{2}}. For submodules $A$, $B$ of a module and a ring element~$z$,
\m{\IndexDvd{A}{B}{z}} is read \emph{the index of $B$ in $A$ divides~$z$}; it is one symbol, defined in
Definition~\ref{d:index}.

\statement{arith}{\phantomsection\label{t:1}\textbf{Theorem 1} (a crude bound for the class number).\enspace
\emph{Let $E$ be a CM field with a CM type, let \m{M\ge1} be an integer whose image in the ring of integers of~$E$
lies in the different of~$E$, and let $\cJ$ be a finite set. Then}
\[\#\,\StdClasses{\Kf{E}{\cJ}}\ \le\ \bigl(1+2\,\sbc{E}{\cJ}{M}^{3}\bigr)^{4\,\numInf{E}\,(\#\cJ)^{2}}.\]

\vspace{2pt}\phantomsection\label{t:2}\textbf{Theorem 2} (the determinant criterion).\enspace
\emph{Let $\cO$ be a valuation domain, $\cJ$ a finite set, $z$ an element of~$\cO$, and $G$ a square matrix over~$\cO$
with rows and columns indexed by~$\cJ$; in words, a matrix~$G$ in \m{\MJ(\cO)}. Let $B$ be a submodule
of~\m{\cO^{\cJ}}. If \m{\det G} divides~$z$ and \m{G(\cO^{\cJ})\subseteq B}, then
\m{\IndexDvd{\cO^{\cJ}}{B}{z}}.}}
\leanline{\lean{cardclassesle}\sep \lean{ofdetdvd}}

In Theorem~1 the number of elements is taken as the formal text takes it, a natural number that would be~$0$ for
an infinite set; the proof shows that the set is finite (§\ref{s:notes}, note~1). Between the two theorems stands
the calculus of §\ref{s:calculus}: inclusion, pull-back, transitivity, push-forward, sums and composition
(Lemmas \ref{l:incl} to~\ref{l:comp}), and the additivity of the \emph{type} of an endomorphism under composition
(Proposition~\ref{p:prod}).

\sect{The setting}{s:setting}
\begin{conventions}\phantomsection\label{c:conventions}
The natural numbers are \m{0,1,2,\dots}. Every ring is a commutative ring with~$1$. In §\ref{s:calculus} and
§\ref{s:det}, $\cO$ is a commutative ring, $V$ and $V'$ are $\cO$-modules, and a submodule is an $\cO$-submodule
of a module. For submodules $A$, $B$ of one module, \m{A+B} is their sum; for an element~$w$ of~$\cO$, \m{wA} is the
submodule of the multiples \m{wv}, \m{v\in A}; for a linear map~$G$, \m{G(A)} is the image and \m{G^{-1}(B)} the
preimage of a submodule. A \kw{valuation domain} is an integral domain in which, of any two elements, one divides
the other. A finite set indexes coordinates: for a finite set~$\cJ$, \m{\cO^{\cJ}} is the module of families
\m{(v_i)_{i\in\cJ}}, \m{\MJ(\cO)} is the ring of matrices with rows and columns indexed by~$\cJ$, acting on
\m{\cO^{\cJ}}, and \m{\GLJ(\cO)} is its group of units. The number of elements of a finite set~$\cS$ is
\m{\#\cS}.

In §\ref{s:class}, $E$ is a CM field: a totally complex quadratic extension of a totally real number field.
Complex conjugation of~$E$ is written \m{v\mapsto\bar v}. A \kw{CM type} of~$E$ is a family \m{\Phi=(\Phi_w)_w} of
field embeddings \m{\Phi_w\colon E\to\CC}, one for each infinite place~$w$ of~$E$, such that $\Phi_w$
induces~$w$. The ring of integers of~$E$ is \m{\cO_E}, its different (over~$\mathbb{Z}$) is the ideal whose
inverse is the dual of~\m{\cO_E} for the trace form. The ring of finite adeles of~$E$ is~$\Af$, and $\Oh$ is the
closure of~\m{\cO_E} in it, the product of the rings of integers of the completions of~$E$ at the finite places;
$E$ is a subring of~$\Af$, and \m{E\cap\Oh=\cO_E}. Complex conjugation extends to~$\Af$ and preserves~$\Oh$. For a
matrix~$h$ over $E$ or over~$\Af$, \m{h^{*}} is its conjugate transpose.
\end{conventions}

\begin{definition}[index divisibility]\label{d:index}
Let $A$ and $B$ be submodules of the $\cO$-module~$V$, and let $z$ be an element of~$\cO$. We write
\[\IndexDvd{A}{B}{z}\]
and say that \kw{the index of $B$ in $A$ divides~$z$}, if for every natural number~$n$ and all elements
\m{y_1,\dots,y_n} of~$A$ there is an \m{n\times n} matrix \m{u=(u_{ii'})} over~$\cO$ such that \m{\det u}
divides~$z$ and
\[\sum_{i'=1}^{n}u_{ii'}\,y_{i'}\in B\qquad\text{for } i=1,\dots,n.\]
\end{definition}
\leanline{\lean{IndexDvd}}

The definition does not ask that $B$ lie in~$A$, nor that either module be free or finitely generated. Since the
combinations \m{\sum_{i'}u_{ii'}y_{i'}} lie in~$A$, the statement \m{\IndexDvd{A}{B}{z}} depends on $B$ only
through \m{A\cap B}; it holds for \m{z=0} (take \m{u=0}), and for every~$z$ when \m{A\subseteq B}
(Lemma~\ref{l:incl}). The bracket alone has no meaning. Remark~\ref{r:dvr} compares the notion with the length
of a quotient in the one case where both are defined. The overview of the handout uses the same pattern for the
index of a family of rows: there, too, a matrix~$u$ with a large determinant moves a family into the integral
forms.

\begin{definition}[type of an endomorphism]\label{d:type}
Let $p$ be an element of~$\cO$, $A$ a submodule of the $\cO$-module~$V$, $\nu$ a natural number, $Y$ an
$\cO$-linear endomorphism of~$V$, and $c$ a real number. We say that
\par\vspace{3pt}\centerline{\HasType{p}{A}{\nu}{Y}{c}}\par\vspace{3pt}\noindent
if \m{Y(A)\subseteq A} and there is a natural number \m{i\le c} with
\m{\IndexDvd{A}{Y(A)+p^{\nu}A}{p^{i}}}.
\end{definition}
\leanline{\lean{HasType}}

When $\cO$ is a discrete valuation ring with uniformiser~$p$ and $A$ is free of finite rank, this says that $Y$
preserves~$A$ and that the cokernel of $Y$ on \m{A/p^{\nu}A} has length at most~$c$ (Remark~\ref{r:dvr}): the
elementary divisors of~$Y$ on~$A$, cut off at~\m{p^{\nu}}, have exponents of sum at most~$c$.

\begin{definition}[the standard definite space and its classes]\label{d:classes}
Let $E$ be a CM field and $\cJ$ a finite set.
\begin{enumerate}[label=\textup{(\alph*)},leftmargin=2.3em,itemsep=2pt]
\item \m{\UJ(\Af)} is the group of the matrices~$h$ in \m{\GLJ(\Af)} with \m{h^{*}h=1}: the finite adelic points
of the unitary group of the hermitian form \m{\sum_{i\in\cJ}\bar v_iv_i} on~\m{E^{\cJ}}. \m{\UJ(E)} is the group of
the matrices~$\gamma$ in \m{\GLJ(E)} with \m{\gamma^{*}\gamma=1}; it is a subgroup of \m{\UJ(\Af)}. The subgroup
\[\Kf{E}{\cJ}\]
of \m{\UJ(\Af)} consists of the elements~$h$ such that $h$ and \m{h^{-1}} have all entries in~$\Oh$; it is the
stabiliser of the standard lattice~\m{\Oh^{\cJ}}.
\item For a subgroup~$K$ of \m{\UJ(\Af)}, the set of \kw{classes at level~$K$} is the set of double cosets
\[\StdClasses{K}=\UJ(E)\backslash\UJ(\Af)/K .\]
\item \m{\numInf{E}} is the number of infinite places of~$E$; this is defined for every number field.
\item For a natural number~$M$, \m{\sbc{E}{\cJ}{M}=6\,M\,(\numInf{E}\cdot\#\cJ)^{2}}, a natural number.
\end{enumerate}
\end{definition}
\leanline{\lean{Kf}\sep \lean{StdClasses}\sep \lean{numInf}\sep \lean{sbc}\sep these four definitions belong to the
background library (units 122, 102, 10 and~99)}

The classes at level \m{\Kf{E}{\cJ}} correspond to the isometry classes of the hermitian lattices in~\m{E^{\cJ}}
that are everywhere locally isometric to~\m{\cO_E^{\cJ}}: the class of~$h$ is that of the lattice
\m{E^{\cJ}\cap h\,\Oh^{\cJ}}. The number of classes is the class number of the genus.

\statement{inksoft}{\textbf{Inputs from other units.} The four results below are taken for granted in this note.

\vspace{4pt}\phantomsection\label{i:1}\textbf{Input 1} (short bases; unit 99, background).\enspace
\emph{Let $E$ be a CM field with a CM type, $\cJ$ a finite set, and \m{M\ge1} an integer that, as an element of
\m{\cO_E}, lies in the different of~$E$. For every $h$ in \m{\UJ(\Af)} there is a matrix~$L$ in \m{\GLJ(E)} such that
\m{h^{-1}L} has all entries in~$\Oh$ and \m{|\varphi(L_{ii'})|\le\sbc{E}{\cJ}{M}} for every ring homomorphism
\m{\varphi\colon E\to\CC} and all \m{i,i'} in~$\cJ$.}

\vspace{3pt}\phantomsection\label{i:2}\textbf{Input 2} (rational approximation of an adele; unit 99,
background).\enspace
\emph{Let $E$ be a CM field with a CM type \m{(\Phi_w)_w}, and let $\delta$ be a nonzero element of the different
of~$E$. For every family \m{(z_w)_w} of complex numbers indexed by the infinite places of~$E$ and every finite
adele~$\eta$ in~$\Af$ there is $\alpha$ in~$E$ with \m{\eta-\alpha\in\Oh} and
\m{|z_w-\Phi_w(\alpha)|\le(2\,\numInf{E})^{2}\,|\Phi_w(\delta)|} for every infinite place~$w$.}

\vspace{3pt}\phantomsection\label{i:3}\textbf{Input 3} (bounded integral vectors are few; unit 148).\enspace
\emph{Let $E'$ be a number field with an embedding \m{\Phi_v\colon E'\to\CC} inducing~$v$ for each infinite
place~$v$. Let $\cS$ be a set, $\cS'$ a subset of it, $\cJ$ a finite set, $\omega$ a real number with
\m{0\le\omega}, and $F$ a map from $\cS$ to \m{E'^{\cJ}}. Suppose that for every $c$ in~$\cS'$ and every $i$
in~$\cJ$ the coordinate \m{F(c)_i} is an algebraic integer with \m{|\Phi_v(F(c)_i)|\le\omega} for every~$v$, and
that $F$ is injective on~$\cS'$. Then $\cS'$ is finite, with at most
\m{(1+2\omega)^{2\,\#\cJ\,\numInf{E'}}} elements.} (The reading of this lemma by program lists, besides these
arguments, variables of its file that the statement does not mention, among them a matrix; by Lean's rules for
such variables they are not arguments of the lemma.)

\vspace{3pt}\phantomsection\label{i:4}\textbf{Input 4} (Hermite normal form; unit 120, background).\enspace
\emph{\textup{(a)} Let $\mathcal{A}$ be a commutative ring and $\mathcal{R}$ a subring such that
\m{(\mathcal{A},\mathcal{R})} is a valuation pair: for any two elements $y$, $y'$ of~$\mathcal{A}$ there is an
idempotent~$\epsilon$ of~$\mathcal{R}$ with
\m{\mathcal{R}y+\mathcal{R}y'\subseteq\mathcal{R}\,(y-\epsilon(y-y'))}. Let $\mathcal{K}$ and $\mathcal{N}$ be
finite sets, \m{\iota\colon\mathcal{K}\to\mathcal{N}} an injective map, $X$ a matrix over~$\mathcal{A}$ with rows
indexed by~$\mathcal{N}$ and columns by~$\mathcal{K}$, and $\omega$ a map from~$\mathcal{K}$ to the natural numbers
\textup{(}weights\textup{)}. Then there are an invertible matrix~$T$ over~$\mathcal{A}$, indexed
by~$\mathcal{N}$, such that $T$ and \m{T^{-1}} have their entries in~$\mathcal{R}$, and a square
matrix~$\Delta$ over~$\mathcal{A}$, indexed by~$\mathcal{K}$, with \m{\Delta_{ll'}=0} whenever
\m{\omega(l')<\omega(l)}, such that
\[X=T_{\cdot,\iota}\,\Delta\qquad\text{and}\qquad (T^{-1})_{\iota,\cdot}\,X=\Delta ,\]
where \m{T_{\cdot,\iota}} consists of the columns \m{\iota(l)} of~$T$ and \m{(T^{-1})_{\iota,\cdot}} of the rows
\m{\iota(l)} of~\m{T^{-1}}, \m{l\in\mathcal{K}}.
\textup{(b)} If $\mathcal{A}$ is a valuation domain, then \m{(\mathcal{A},\mathcal{A})} is a valuation pair.}
(The reading of~(a) by program also lists a ring with involution, from an earlier part of its file; the statement
does not involve it.)}
\leanline{\lean{shortbasis}\sep \lean{ratnear}\sep \lean{enc}\sep \lean{hermite}\sep \lean{valtop}}

Inputs 1 and~2 are lemmas of the reduction theory that the artifact develops for the standard space: Input~1
comes from a transference argument (a basis of short vectors, from the covering radius of a lattice against the
first minimum of its dual), Input~2 from a fundamental domain for~\m{\cO_E} scaled by the different. Input~3 is a
packing count: distinct algebraic integers are at distance at least~$1$ at some infinite place. Input~4(a), for a
discrete valuation ring and \m{\mathcal{K}=\mathcal{N}}, is the Hermite normal form obtained by column operations;
the valuation pair is the form in which the artifact proves it once for a valuation domain and for $\Oh$
inside~$\Af$. The catalogue of quoted results of this write-up lists Inputs 1, 2 and~4 as elementary lemmas, with
no confirmed place in the literature: \emph{reference to be confirmed}.

\sect{The class number}{s:class}
\begin{proof}[Proof of Theorem 1]
Write \m{\mathrm{Cl}} for the set \m{\StdClasses{\Kf{E}{\cJ}}}, and choose for every class~$\fc$ in it a
representative \m{h_{\fc}} in \m{\UJ(\Af)}.

\emph{Short bases.} By Input~1 there is, for every class~$\fc$, a matrix \m{L_{\fc}} in \m{\GLJ(E)} such that
\[Q_{\fc}=h_{\fc}^{-1}L_{\fc}\ \text{ has its entries in } \Oh\]
and
\[|\varphi((L_{\fc})_{ii'})|\le\sbc{E}{\cJ}{M}\quad\text{for every ring homomorphism } \varphi\colon E\to\CC \text{ and all } i,i'.\]
Let \m{H_{\fc}=L_{\fc}^{*}L_{\fc}}, the Gram matrix of the columns of~\m{L_{\fc}}; it lies in \m{\GLJ(E)}.

\emph{An approximate inverse.} Apply Input~2 with \m{\delta=M} and \m{z_w=0}: for every finite adele~$\eta$ there
is \m{\varrho(\eta)} in~$E$ with
\[\eta-\varrho(\eta)\in\Oh,\qquad|\Phi_w(\varrho(\eta))|\le(2\,\numInf{E})^{2}M\ \text{ for every infinite place } w.\]
Let \m{Z_{\fc}} be the matrix over~$E$ obtained by applying $\varrho$ to each entry of \m{Q_{\fc}^{-1}}. Then
\m{Q_{\fc}^{-1}-Z_{\fc}} has its entries in~$\Oh$, and so, for all classes $\fc$ and~$\fc'$,
\begin{equation}\label{e:int}
Q_{\fc'}Q_{\fc}^{-1}-Q_{\fc'}Z_{\fc}=Q_{\fc'}\bigl(Q_{\fc}^{-1}-Z_{\fc}\bigr)\quad\text{has its entries in } \Oh .
\end{equation}

\emph{The map.} Send a class~$\fc$ to the pair of matrices \m{(H_{\fc},\,H_{\fc}Z_{\fc})}, a vector with
\m{2(\#\cJ)^{2}} coordinates in~$E$. We show that its coordinates are algebraic integers, that they are bounded
by \m{\sbc{E}{\cJ}{M}^{3}} at every infinite place, and that the map is injective. Input~3, for the field~$E$ with
its CM type, the set \m{\cS=\cS'=\mathrm{Cl}}, an index set of \m{2(\#\cJ)^{2}} elements and
\m{\omega=\sbc{E}{\cJ}{M}^{3}}, then shows that \m{\mathrm{Cl}} is finite with at most
\m{(1+2\,\sbc{E}{\cJ}{M}^{3})^{2\cdot2(\#\cJ)^{2}\cdot\numInf{E}}} elements, which is the bound of the theorem.

\emph{Integrality.} Since \m{h_{\fc}^{-1}} is unitary,
\m{H_{\fc}=(h_{\fc}^{-1}L_{\fc})^{*}(h_{\fc}^{-1}L_{\fc})=Q_{\fc}^{*}Q_{\fc}}, and the entries of \m{Q_{\fc}^{*}}
lie in~$\Oh$ because conjugation preserves~$\Oh$. So $H_{\fc}$ has its entries in \m{E\cap\Oh=\cO_E}. In the same
way \m{H_{\fc}Z_{\fc}=Q_{\fc}^{*}\,(Q_{\fc}Z_{\fc})}, and \m{Q_{\fc}Z_{\fc}=1-Q_{\fc}(Q_{\fc}^{-1}-Z_{\fc})} has its
entries in~$\Oh$ by \eqref{e:int} with \m{\fc'=\fc}.

\emph{Size.} Let $w$ be an infinite place. The composite of $\Phi_w$ with complex conjugation of~$E$ is again a
ring homomorphism from $E$ to~$\CC$, so every entry of $L_{\fc}$ and of $L_{\fc}^{*}$ has absolute value at most
\m{\sbc{E}{\cJ}{M}} under~$\Phi_w$. Hence
\[|\Phi_w((H_{\fc})_{ii'})|\le\#\cJ\cdot\sbc{E}{\cJ}{M}^{2},\qquad
|\Phi_w((H_{\fc}Z_{\fc})_{ii'})|\le(\#\cJ)^{2}\,\sbc{E}{\cJ}{M}^{2}\,(2\,\numInf{E})^{2}M .\]
Here \m{\#\cJ\ge1}, as there is an entry, and \m{M\ge1}, \m{\numInf{E}\ge1}. So
\[\#\cJ\le\sbc{E}{\cJ}{M}\qquad\text{and}\qquad 4M\,\numInf{E}^{2}(\#\cJ)^{2}\le 6M(\numInf{E}\cdot\#\cJ)^{2}=\sbc{E}{\cJ}{M},\]
and both bounds are at most the cube of \m{\sbc{E}{\cJ}{M}}.

\emph{Injectivity.} Let $\fc$ and $\fc'$ be classes with \m{H_{\fc}=H_{\fc'}} and
\m{H_{\fc}Z_{\fc}=H_{\fc'}Z_{\fc'}}. As \m{H_{\fc}} is invertible, \m{Z_{\fc}=Z_{\fc'}}. Put
\m{\gamma=L_{\fc}L_{\fc'}^{-1}}, a matrix in \m{\GLJ(E)}. Then
\[\gamma^{*}\gamma=(L_{\fc'}^{-1})^{*}\,H_{\fc}\,L_{\fc'}^{-1}=(L_{\fc'}^{-1})^{*}\,H_{\fc'}\,L_{\fc'}^{-1}=1,\]
so $\gamma$ lies in \m{\UJ(E)}. Put \m{\kappa=(\gamma h_{\fc'})^{-1}h_{\fc}}, an element of \m{\UJ(\Af)}. Then
\[\kappa=h_{\fc'}^{-1}L_{\fc'}L_{\fc}^{-1}h_{\fc}=Q_{\fc'}Q_{\fc}^{-1},\]
and by \eqref{e:int}, used for the pairs \m{(\fc,\fc')} and \m{(\fc',\fc')}, and by \m{Z_{\fc}=Z_{\fc'}},
\[\kappa=1+\bigl(Q_{\fc'}Q_{\fc}^{-1}-Q_{\fc'}Z_{\fc}\bigr)-\bigl(Q_{\fc'}Q_{\fc'}^{-1}-Q_{\fc'}Z_{\fc'}\bigr)\]
has its entries in~$\Oh$. The inverse of the unitary matrix~$\kappa$ is \m{\kappa^{*}}, which has its entries
in~$\Oh$ as well, so $\kappa$ lies in \m{\Kf{E}{\cJ}}. From \m{h_{\fc}=\gamma\,h_{\fc'}\,\kappa} the two
representatives lie in one double coset, and \m{\fc=\fc'}.
\end{proof}
\leanline{\lean{cardclassesle}\sep the formal proof has these five parts in this order}

The bound is far from the truth: the mass formula gives the order of magnitude of the class number. What the
argument of the write-up needs is a bound that is explicit and at most exponential in a polynomial of
\m{\numInf{E}}, \m{\#\cJ} and~$M$.

\sect{The calculus of index divisibility}{s:calculus}
In this section $\cO$ is a commutative ring and $V$, $V'$ are $\cO$-modules, as in the conventions.

\begin{lemma}[inclusion]\label{l:incl}
Let $A$ and $B$ be submodules of~$V$ and $z$ an element of~$\cO$. If \m{A\subseteq B}, then
\m{\IndexDvd{A}{B}{z}}.
\end{lemma}
\begin{proof}
For \m{y_1,\dots,y_n} in~$A$ the identity matrix serves: its determinant~$1$ divides~$z$, and each $y_i$ lies
in~$B$.
\end{proof}
\leanline{\lean{ofle}}

\begin{lemma}[pull-back]\label{l:comap}
Let $A$, $B$ be submodules of~$V$, let $A'$, $B'$ be submodules of~$V'$, let $z$ be an element of~$\cO$, and let
\m{G\colon V'\to V} be an $\cO$-linear map. Suppose that \m{\IndexDvd{A}{B}{z}}, that \m{G(A')\subseteq A}, and
that every $v$ in~$A'$ with \m{G(v)\in B} lies in~$B'$. Then \m{\IndexDvd{A'}{B'}{z}}.
\end{lemma}
\begin{proof}
Let \m{y_1,\dots,y_n} lie in~$A'$. The elements \m{G(y_i)} lie in~$A$, so there is a matrix~$u$ with \m{\det u}
dividing~$z$ and \m{\sum_{i'}u_{ii'}G(y_{i'})\in B} for all~$i$. The element \m{v_i=\sum_{i'}u_{ii'}y_{i'}} lies
in~$A'$ and \m{G(v_i)\in B}; hence \m{v_i\in B'}.
\end{proof}
\leanline{\lean{comap}}

\begin{lemma}[transitivity]\label{l:trans}
Let $A$, $B$, $C$ be submodules of~$V$ and let $z$ and $w$ be elements of~$\cO$. If \m{\IndexDvd{A}{B}{z}} and
\m{\IndexDvd{B}{C}{w}}, then \m{\IndexDvd{A}{C}{zw}}.
\end{lemma}
\begin{proof}
Let \m{y_1,\dots,y_n} lie in~$A$. There is $u$ with \m{\det u} dividing~$z$ such that the elements
\m{y'_i=\sum_{i'}u_{ii'}y_{i'}} lie in~$B$; and for these there is $u'$ with \m{\det u'} dividing~$w$ such that
\m{\sum_{i}u'_{li}y'_{i}\in C} for all~$l$. The matrix \m{u'u} has determinant \m{\det u'\cdot\det u}, which
divides~\m{zw}, and \m{\sum_{i'}(u'u)_{li'}y_{i'}=\sum_iu'_{li}y'_i} lies in~$C$.
\end{proof}
\leanline{\lean{trans}}

\begin{lemma}[push-forward]\label{l:map}
Let $A$, $B$ be submodules of~$V$, let $A'$, $B'$ be submodules of~$V'$, let $z$ be an element of~$\cO$, and let
\m{G\colon V\to V'} be an $\cO$-linear map. Suppose that \m{\IndexDvd{A}{B}{z}}, that
\m{A'\subseteq G(A)+B'}, and that \m{G(B)\subseteq B'}. Then \m{\IndexDvd{A'}{B'}{z}}.
\end{lemma}
\begin{proof}
Let \m{y_1,\dots,y_n} lie in~$A'$, and write \m{y_i=G(v_i)+b_i} with $v_i$ in~$A$ and $b_i$ in~$B'$. There is $u$
with \m{\det u} dividing~$z$ and \m{\sum_{i'}u_{ii'}v_{i'}\in B} for all~$i$. Then
\[\sum_{i'}u_{ii'}y_{i'}=G\Bigl(\sum_{i'}u_{ii'}v_{i'}\Bigr)+\sum_{i'}u_{ii'}b_{i'}\ \in\ G(B)+B'\subseteq B'.\qedhere\]
\end{proof}
\leanline{\lean{map}}

\begin{lemma}[finite sums]\label{l:sup}
Let $\cJ$ be a set, let \m{G_i\colon V\to V'}, for $i$ in~$\cJ$, be $\cO$-linear maps, let $A$ be a submodule
of~$V$, $B'$ a submodule of~$V'$ and $z$ an element of~$\cO$. Suppose that
\m{\IndexDvd{A}{G_i^{-1}(B')}{z}} for every $i$ in~$\cJ$. Then for every finite subset $\cS$ of~$\cJ$,
\[\IndexDvd{\textstyle\sum_{i\in\cS}G_i(A)+B'}{B'}{z^{\#\cS}} .\]
\end{lemma}
\begin{proof}
By induction on~$\cS$. For the empty set the left module is~$B'$, and Lemma~\ref{l:incl} applies. Let
\m{\cS=\cS_0\cup\{i\}} with $i$ not in~$\cS_0$, and put \m{A'_0=\sum_{l\in\cS_0}G_l(A)+B'} and
\m{A'=G_i(A)+A'_0}. Lemma~\ref{l:map}, for the map~$G_i$ and the submodules $A$, \m{G_i^{-1}(B')} of~$V$ and $A'$,
$A'_0$ of~$V'$, gives \m{\IndexDvd{A'}{A'_0}{z}}: indeed \m{A'\subseteq G_i(A)+A'_0} and
\m{G_i(G_i^{-1}(B'))\subseteq B'\subseteq A'_0}. By the induction hypothesis
\m{\IndexDvd{A'_0}{B'}{z^{\#\cS_0}}}, and Lemma~\ref{l:trans} gives the claim with
\m{z\cdot z^{\#\cS_0}=z^{\#\cS}}.
\end{proof}
\leanline{\lean{sup}\sep the formal statement writes the left module as the span of the union of the images
\m{G_i(A)}, \m{i\in\cS}, and of~$B'$}

\begin{lemma}[composition]\label{l:comp}
Let $A$ and $B$ be submodules of~$V$, let $z$, $z'$ and $w$ be elements of~$\cO$, and let \m{G\colon V\to V} be an
$\cO$-linear map with \m{G(A)\subseteq A}. If \m{\IndexDvd{A}{G(A)+wA}{z}} and \m{\IndexDvd{A}{B+wA}{z'}}, then
\m{\IndexDvd{A}{G(B)+wA}{zz'}}.
\end{lemma}
\begin{proof}
Apply Lemma~\ref{l:map} to the second hypothesis and the map~$G$, with \m{A'=G(A)+wA} and \m{B'=G(B)+wA}. The
inclusion \m{A'\subseteq G(A)+B'} is clear, and
\[G(B+wA)=G(B)+w\,G(A)\subseteq G(B)+wA\]
because \m{G(A)\subseteq A}. This gives \m{\IndexDvd{G(A)+wA}{G(B)+wA}{z'}}. Lemma~\ref{l:trans}, with the first
hypothesis, gives the claim.
\end{proof}
\leanline{\lean{comp}}

\begin{proposition}[types add under composition]\label{p:prod}
Let $\cJ$ be a set, and for every $i$ in~$\cJ$ let $Y_i$ be an $\cO$-linear endomorphism of~$V$ and $c_i$ a real
number. Let $p$ be an element of~$\cO$ \textup{(}the formal text calls it~$z$; we write $p$ as in
Definition~\ref{d:type}, and \m{z=p} here\textup{)}, $A$ a submodule of~$V$ and $\nu$ a natural number. Let
\m{(i_1,\dots,i_e)} be a finite list of elements of~$\cJ$, with \m{e\ge0}; repetitions are allowed. Suppose that
\HasType{p}{A}{\nu}{Y_{i_l}}{c_{i_l}} for \m{l=1,\dots,e}. Then
\HasType{p}{A}{\nu}{Y_{i_1}\circ\dots\circ Y_{i_e}}{c_{i_1}+\dots+c_{i_e}}; for \m{e=0} the composite is the
identity and the sum is~$0$.
\end{proposition}
\begin{proof}
By induction on~$e$. For \m{e=0}: the identity preserves~$A$, and \m{\IndexDvd{A}{A+p^{\nu}A}{p^{0}}} by
Lemma~\ref{l:incl}, with the natural number \m{0\le0}. For the step put \m{Y=Y_{i_1}} and
\m{Y'=Y_{i_2}\circ\dots\circ Y_{i_e}}. By the hypothesis on~$Y$ there is a natural number \m{i\le c_{i_1}} with
\m{\IndexDvd{A}{Y(A)+p^{\nu}A}{p^{i}}}, and \m{Y(A)\subseteq A}; by the induction hypothesis \m{Y'(A)\subseteq A} and
there is a natural number \m{i'\le c_{i_2}+\dots+c_{i_e}} with \m{\IndexDvd{A}{Y'(A)+p^{\nu}A}{p^{i'}}}. Then
\m{Y(Y'(A))\subseteq Y(A)\subseteq A}, and Lemma~\ref{l:comp}, with \m{G=Y}, \m{B=Y'(A)} and \m{w=p^{\nu}}, gives
\m{\IndexDvd{A}{Y(Y'(A))+p^{\nu}A}{p^{i+i'}}}, where \m{i+i'\le c_{i_1}+\dots+c_{i_e}}.
\end{proof}
\leanline{\lean{prod}}

\sect{The determinant criterion}{s:det}
\begin{proof}[Proof of Theorem 2]
Let \m{y_1,\dots,y_n} be elements of~\m{\cO^{\cJ}}; write \m{(y_i)_l} for the coordinates, \m{l\in\cJ}. Let
$\mathcal{N}$ be the disjoint union of \m{\{1,\dots,n\}} and~$\cJ$, and let $X$ be the square matrix over~$\cO$,
indexed by~$\mathcal{N}$, with the blocks
\[X=\begin{pmatrix}1&-(y_i)_l\\[1pt]0&G^{\mathsf{T}}\end{pmatrix},\]
where $1$ is the identity matrix of size~$n$, the upper right block has the entry \m{-(y_i)_l} in row~$i$ and
column~$l$, and \m{G^{\mathsf{T}}} is the transpose of~$G$. By Input~4(b) the pair \m{(\cO,\cO)} is a valuation
pair. Input~4(a), for \m{\mathcal{K}=\mathcal{N}}, the identity map~$\iota$, the matrix~$X$ and the weights
\m{\omega=1} on \m{\{1,\dots,n\}} and \m{\omega=0} on~$\cJ$, gives an invertible matrix~$T$ over~$\cO$ whose
inverse has entries in~$\cO$, and a matrix~$\Delta$ with
\[X=T\Delta,\qquad T^{-1}X=\Delta,\qquad \Delta_{il}=0\ \text{ for } i\in\{1,\dots,n\},\ l\in\cJ .\]
Write \m{\Delta=\bigl(\begin{smallmatrix}\Delta_{11}&0\\\Delta_{21}&\Delta_{22}\end{smallmatrix}\bigr)} and
\m{T^{-1}=\bigl(\begin{smallmatrix}T'_{11}&T'_{12}\\T'_{21}&T'_{22}\end{smallmatrix}\bigr)} in blocks. We claim that
\m{u=\Delta_{11}} is a matrix as Definition~\ref{d:index} asks for.

\emph{The determinant.} As $X$ is block upper triangular, \m{\det X=\det G^{\mathsf{T}}=\det G}; as $\Delta$ is
block lower triangular, \m{\det X=\det T\cdot\det\Delta_{11}\cdot\det\Delta_{22}}. So \m{\det u} divides
\m{\det G}, which divides~$z$.

\emph{The combinations.} The first block row of \m{T^{-1}X=\Delta} reads \m{T'_{11}=\Delta_{11}} in the first block
column and \m{-T'_{11}\,(y_i)_l+T'_{12}\,G^{\mathsf{T}}=0} in the second. Hence, for \m{i=1,\dots,n} and
\m{l\in\cJ},
\[\sum_{i'=1}^{n}u_{ii'}\,(y_{i'})_l=\sum_{l'\in\cJ}(T'_{12})_{il'}\,G_{ll'} ,\]
that is, \m{\sum_{i'}u_{ii'}y_{i'}=G(t_i)} for the vector \m{t_i=((T'_{12})_{il'})_{l'}} of~\m{\cO^{\cJ}}. It
lies in \m{G(\cO^{\cJ})\subseteq B}.
\end{proof}
\leanline{\lean{ofdetdvd}\sep \lean{hermite}\sep \lean{valtop}}

\begin{remark}[comparison with the length of a quotient; not in the formal text]\label{r:dvr}
Let $\cO$ be a discrete valuation ring with uniformiser~$p$, let $A$ be a free $\cO$-module of finite rank~$n$ and
$B$ a submodule of~$A$ with \m{A/B} of finite length~$e$. Then, for $z$ in~$\cO$: \m{\IndexDvd{A}{B}{z}} holds if and
only if \m{p^{e}} divides~$z$. Indeed, $B$ is free of rank~$n$, so in a basis of~$A$ it is the image of a
matrix~$G$ with \m{\det G} a unit times~\m{p^{e}}, and Theorem~2 gives one direction. Conversely, apply
Definition~\ref{d:index} to a basis \m{y_1,\dots,y_n} of~$A$: if \m{z\ne0}, the matrix~$u$ defines an injective
endomorphism of~$A$ with image in~$B$ and cokernel of length the valuation of \m{\det u}, so $e$ is at most the
valuation of~$z$; and \m{p^{e}} divides~$0$.
\end{remark}

\sect{Where the results are used}{s:use}
Theorem~1 is used once, in the unit of the columns (unit~145), in a bound for the number of classes of
standard definite spaces at the levels that contain a principal level: such a level has finite index in
\m{\Kf{E}{\cJ}}, and the number of classes is at most that index times the bound of Theorem~1.

The index language is used in the units that prove Theorem~A of the write-up. The determinant criterion is
applied there with diagonal matrices (units 150 and~141); pull-back and push-forward transport index statements
along the maps between spaces of forms (units 141 and~145); the lemma on finite sums is the step, in the unit
of the endgame (unit~151), from one translate of a Poincar\'e series to finitely many, which is where the
number of translates enters the index bound of Theorem~A(b); and Proposition~\ref{p:prod} adds the types of the
operators attached to the marked primes (unit~145). Unit~144, on the elementary divisors of a Frobenius-twisted
Hecke operator, uses Theorem~2, transitivity, push-forward and composition, and states its result with
Definition~\ref{d:type}.

\sect{Notes on the formal text}{s:notes}
\begin{enumerate}[label=\textup{\arabic*.},leftmargin=1.6em,itemsep=5pt]
\item \emph{The number of classes.} The formal statement of Theorem~1 bounds the cardinality of the set of
classes as a natural number, which the library defines to be~$0$ for an infinite set. The formal proof bounds
the cardinality as an extended natural number, through Input~3, and so proves that the set is finite; the
statement does not say it.
\item \emph{Letters.} The formal text calls the integer of Theorem~1 $N$ and the index set~$n$; here they are $M$
and~$\cJ$, the letter $N$ being kept in this write-up for the level. The submodules $M$, $N$, $P$, $M'$, $N'$ of
the formal text are $A$, $B$, $C$, $A'$, $B'$ here; its modules $V$ and $U$ are $V$ and $V'$; its scalar~$a$
is~$w$; its linear maps~$f$ are~$G$. In Proposition~\ref{p:prod} the formal text writes $z$, $s$, $T$, $t$, $L$ for
what is here $p$, $\nu$, \m{(Y_i)}, \m{(c_i)} and the list.
\item \emph{The arguments of the index relation.} The file declares $M$, $N$, $z$ as implicit variables and then
makes them explicit for this one definition. A reading of the text by program therefore lists each of the three
twice. Lean takes each once, as an explicit argument: every use in the artifact gives exactly three arguments.
This was not confirmed by asking Lean.
\item \emph{Inferred types.} In the definition of the type the formal text does not write the types of $p$, of the
submodule and of the exponent; Lean infers them. The uses (Proposition~\ref{p:prod}, where $p$ is an element
of the ring) show that $p$ is a ring element and the exponent a natural number.
\item \emph{Families of length zero.} Definition~\ref{d:index} quantifies over all~$n$, including \m{n=0}, where
there is nothing to check. The families are indexed by \m{\{0,\dots,n-1\}} in the formal text.
\item \emph{Lemma~\ref{l:sup}.} Its hypothesis is asked for every $i$ in~$\cJ$, not only for those in~$\cS$, and
$\cJ$ is an arbitrary set. The formal proof is one line of induction that calls push-forward and transitivity,
as here.
\item \emph{The injectivity step of Theorem~1.} The formal proof gives the rational unitary matrix and the
integral adelic matrix as terms and leaves it to Lean to see that they relate the two representatives. The
arrangement written above, \m{h_{\fc}=\gamma h_{\fc'}\kappa} with \m{\gamma=L_{\fc}L_{\fc'}^{-1}} and
\m{\kappa=Q_{\fc'}Q_{\fc}^{-1}}, is read off these terms: the unitarity is proved for the product
\m{L_{\fc}L_{\fc'}^{-1}}, and the integrality for the combination of the two instances of \eqref{e:int} displayed
in the proof, which is \m{Q_{\fc'}Q_{\fc}^{-1}}. The library's lemma on double cosets, which fixes the order, was
not opened: the library is not on the machine where this note was written.
\item \emph{The size estimates.} The formal proof closes the two inequalities of the paragraph \emph{Size} by
automation. The intermediate bounds written there are ours; they were checked by hand.
\item \emph{What an assembly assembles.} Lemma~\ref{l:comp} is push-forward followed by transitivity, and
Lemma~\ref{l:sup} is an induction on these two; their content lies in Lemmas \ref{l:trans} and~\ref{l:map}.
\item \emph{The place of Theorem~1.} The header of its file counts it with the reduction theory of the
background (Mahler and Godement compactness, finite covolume, adelic Minkowski theorems), a known subject; it
stands in this unit because its proof needs the count of Input~3, which belongs to the argument.
\end{enumerate}

\subhead{How this note was written.} From the dossier of the unit: its eleven statement readings, its five
structured proofs and the reader's list of doubts. The program check of the dossier covered \UnitSteps{} leaf steps
of \UnitProofs{} proofs. A second reader compared the eleven readings with the Lean text and sampled the proofs;
the remarks of that reading are taken up in the paragraph after Definition~\ref{d:index} and in the hypotheses
of Theorem~2. I read both files of the unit completely, and in the background library the statements and
definitions that the unit rests on: the unitary groups, the level subgroup, the rational points and the classes,
the constant and the two lemmas of Inputs 1 and~2, the count of Input~3 with the variables of its file, the
valuation pairs and the Hermite form. I did not read the proofs of the Inputs. The three lemmas without a
structured proof (pull-back, inclusion, composition) and Lemma~\ref{l:sup} were read in the source. The proofs
above follow the formal ones; the proof of Theorem~2 writes out in blocks what the formal proof obtains by
simplification, and Remark~\ref{r:dvr} is mine. Lean was not run for this note.

\head{Appendix A. Concordance with the Lean artifact}
{\fontsize{8.6}{11.6}\selectfont\renewcommand{\leanfont}{\ttfamily\fontsize{7.7}{10}\selectfont\color{leancol}}
\setlength{\LTleft}{0pt}\setlength{\LTright}{\fill}\setlength{\tabcolsep}{0pt}
\begin{longtable}{@{}p{30mm}p{56mm}p{43mm}r@{}}
\textbf{In this note} & \textbf{Lean declaration} & \textbf{File} & \textbf{Lines}\\\midrule\endhead
\ConcordanceRows
\end{longtable}}

\subhead{The two theorems and the two definitions as they are written in the artifact.}
\quoted{cardclassesle}
\quoted{ofdetdvd}
\quoted{IndexDvd}
\quoted{HasType}
In these texts {\leanfont ⊤} is the whole module, {\leanfont ⊔} the sum of submodules, {\leanfont M.map y} the
image and {\leanfont p \textasciicircum{} n • M} the multiples; {\leanfont Fin n} is the set \m{\{0,\dots,n-1\}}.

\head{Notation}
\begingroup\fontsize{8.6}{10.9}\selectfont\raggedright
\setlength{\columnsep}{9mm}\setlength{\parskip}{0pt}
\begin{multicols}{2}
\nt{$\cO$}{a commutative ring; in Theorem~2 a valuation domain; $\mathcal{O}$ with a subscript is a ring of integers}
\nt{$V$, $V'$}{$\cO$-modules}
\nt{$A$, $B$, $C$, $A'$, $B'$, $A'_0$}{submodules}
\nt{$z$, $z'$, $w$}{elements of~$\cO$}
\nt{\m{\IndexDvd{A}{B}{z}}}{the index of $B$ in $A$ divides~$z$ (Definition~\ref{d:index})}
\nt{$n$; \m{y_1,\dots,y_n}; $y'_i$, $v_i$, $b_i$, $t_i$}{the length of a family; a family of elements of a module; further elements}
\nt{$u$, $u'$, \m{u_{ii'}}}{square matrices over~$\cO$ and their entries}
\nt{$i$, $i'$, $l$, $l'$}{indices; also natural numbers bounding exponents}
\nt{$G$, $G_i$, \m{G^{\mathsf{T}}}}{a linear map, also a square matrix; a family of linear maps; the transpose}
\nt{$p$, $\nu$, $Y$, $c$}{a ring element, a natural number, an endomorphism, a real number (Definition~\ref{d:type})}
\nt{$Y_i$, $c_i$, $Y'$, $e$}{endomorphisms, real numbers; a composite; the length of a list, in Remark~\ref{r:dvr} a length}
\nt{$\cJ$, $\cS$, $\cS'$, $\cS_0$, \m{\#\cS}}{finite sets or sets of indices; the number of elements}
\nt{\m{\cO^{\cJ}}, \m{\MJ(\cO)}, \m{\GLJ(\cO)}}{families, matrices and invertible matrices indexed by~$\cJ$}
\nt{$E$, $\cO_E$, $\Af$, $\Oh$}{a CM field, its integers, its finite adeles, the integral adeles}
\nt{$\Phi_w$, $\varphi$, $w$, $v$}{the embeddings of a CM type; a ring homomorphism to~$\CC$; infinite places ($w$ is also a ring element in §\ref{s:calculus})}
\nt{$M$}{an integer \m{\ge1} in the different of~$E$}
\nt{\m{\UJ(\Af)}, \m{\UJ(E)}, \m{\Kf{E}{\cJ}}}{the unitary group of the standard form: adelic points, rational points, the stabiliser of the standard lattice}
\nt{\m{\StdClasses{K}}, $K$}{the classes at level~$K$; a subgroup of \m{\UJ(\Af)}}
\nt{\m{\numInf{E}}}{the number of infinite places of~$E$}
\nt{\m{\sbc{E}{\cJ}{M}}}{the constant \m{6M(\numInf{E}\cdot\#\cJ)^{2}}}
\nt{$\fc$, $\fc'$, $h_{\fc}$, $h$, \m{h^{*}}}{classes; a representative; an adelic matrix; the conjugate transpose}
\nt{$L_{\fc}$, $L$, $Q_{\fc}$, $H_{\fc}$, $Z_{\fc}$}{a short basis; \m{h_{\fc}^{-1}L_{\fc}}; the Gram matrix; the approximate inverse}
\nt{$\varrho$, $\eta$, $\alpha$, $\delta$, $z_w$}{the approximation map; an adele; an element of~$E$; an element of the different; complex numbers}
\nt{$\gamma$, $\kappa$}{a rational unitary matrix; an element of the level subgroup}
\nt{$E'$, $\Phi_v$, $F$, $\omega$}{in Input~3: a number field, its embeddings, a map, a bound; in Input~4, $\omega$ are weights}
\nt{$\mathcal{A}$, $\mathcal{R}$, $\mathcal{K}$, $\mathcal{N}$, $\iota$, $\epsilon$}{in Input~4: a ring, a subring, finite sets, an injection, an idempotent}
\nt{$X$, $T$, $\Delta$, $T'_{11}$, $\Delta_{11}$}{the matrix to be reduced, the transformation, the triangular form, their blocks}
\nt{$y$, $y'$, $\bar v$}{in Input~4 two ring elements; the complex conjugate}
\nt{$N$, $P$, $U$, $a$, $f$, $s$, $t$, $L$}{letters of the formal text, used only in §\ref{s:notes}, note~2}
\end{multicols}
\endgroup

\end{document}
